% ECONOMICS_PROBLEM: {"id": "O33-478", "title": "Technology diffusion and intangible investment", "jel_code": "O33", "source_id": "OP-0264"}
% FIRST_POSTED: 2026-10-09
% ATTEMPT_STATUS: unresolved
% ENTRY_KIND: research_attempt
% !TEX program = pdflatex
\documentclass[11pt]{article}
\usepackage[T1]{fontenc}
\usepackage[utf8]{inputenc}
\usepackage[margin=26mm]{geometry}
\usepackage{amsmath,amssymb,amsthm,mathtools}
\usepackage[colorlinks=true,linkcolor=blue,citecolor=blue,urlcolor=blue]{hyperref}
\allowdisplaybreaks
\newtheorem{theorem}{Theorem}
\newtheorem{proposition}[theorem]{Proposition}
\newtheorem{lemma}[theorem]{Lemma}
\newtheorem{corollary}[theorem]{Corollary}
\theoremstyle{remark}
\newtheorem{remark}[theorem]{Remark}
\newcommand{\E}{\mathbb E}
\newcommand{\R}{\mathbb R}
\newcommand{\Ind}{\mathbf 1}
\newcommand{\Var}{\operatorname{Var}}
\newcommand{\supp}{\operatorname{supp}}
\DeclareMathOperator*{\argmax}{arg\,max}
\title{Feasible acceleration and the accounting of complementary intangible capital}
\author{GPT 6 Astra Ultra}
\date{9 October 2026}
\begin{document}
\maketitle
\noindent\textbf{Problem ID:} \texttt{O33-478}. \textbf{Status:} Unresolved; restricted results with complete proofs.

\section{Question and separation of targets}
The target separates the true-productivity difference between feasible accelerated and actual adoption, $D_H$, from the accounting discrepancy on the actual path, $M_H$. We derive an exact discrepancy formula, a complementary-capital reachability bound, and a resource-feasible example in which accelerated adoption lowers a finite-horizon productivity residual. We also show that actual-path accounting cannot identify the accelerated counterfactual.

Brynjolfsson, Korinek, and Agrawal ask how adoption and diffusion influence growth and which policies promote optimal diffusion \cite[Section 3.1, printed p. 10]{ai}. That paragraph was inspected. It motivates the diffusion question; the intangible-accounting model and proofs here are explicitly additional constructions.

\section{An accounting convention with both sides recorded}
At dates $t=0,\ldots,H$, let $C_t>0$ be final marketed output excluding own-account intangible formation, and $J_t\geq0$ the latter's real value in the same fixed-price units. Full output is $Y_t=C_t+J_t$. Intangible stock satisfies
\[
 S_{t+1}=(1-\delta)S_t+J_t,\qquad S_0>0,\quad0\leq\delta\leq1.
\]
Only paths with $S_H>0$ are used for the horizon-$H$ residual. This stock recursion requires investment and stock to be valued in compatible units. It does not add $J_t$ again to final consumption. Let $K_H,L_H>0$ be observed physical capital and labor, let $S_{\rm ref}>0$ be a fixed stock normalization, and fix input exponents $\alpha,\nu,\gamma\geq0$. Define
\[
 A_H=\frac{C_H+J_H}{K_H^\alpha L_H^\nu(S_H/S_{\rm ref})^\gamma},
 \qquad
 \widetilde A_H=\frac{C_H}{K_H^\alpha L_H^\nu}.
\]
Measured accounts here expense all of $J_H$ and omit the intangible stock from the input index. These are declared residuals, not an assertion that both are the same structural technology parameter. Their ratio is meaningful only with the stated units and normalization.

\begin{proposition}[Exact accounting discrepancy]\label{accounting}
On a fixed real allocation path,
\[
 M_H=\log A_H-\log\widetilde A_H
    =\log(1+J_H/C_H)-\gamma\log(S_H/S_{\rm ref}).
\]
It is nonnegative exactly when
$1+J_H/C_H\geq(S_H/S_{\rm ref})^\gamma$. Thus omission of intangible investment does not have an unconditional sign for the productivity discrepancy.
\end{proposition}
\begin{proof}
Divide the two residual definitions. The physical capital and labor factors cancel, leaving $(1+J_H/C_H)(S_H/S_{\rm ref})^{-\gamma}$. Taking logarithms gives the identity. Positivity of all logged factors and monotonicity of logarithms yield the sign equivalence.
\end{proof}

For example, with $C_H=1$, $J_H=0$, $S_H/S_{\rm ref}=2$, and $\gamma=1/2$, $M_H=-\tfrac12\log2$. With $C_H=1$, $J_H=1$, $S_H=S_{\rm ref}$, it is $\log2$. These are accounting configurations; the next section imposes actual investment constraints. The numerator omission lowers measured output, while the denominator omission raises its residual relative to a full-input calculation. Both must be recorded before attributing a measured lag.

\begin{proposition}[Linked-history error bounds]
Suppose $C_H$ is known, $\delta$ and $S_0$ are known, and each investment satisfies $J_t\in[\ell_t,u_t]$, $0\leq\ell_t\leq u_t<\infty$. Define
\[
 S^-=(1-\delta)^HS_0+\sum_{t=0}^{H-1}(1-\delta)^{H-1-t}\ell_t,
 \quad
 S^+=(1-\delta)^HS_0+\sum_{t=0}^{H-1}(1-\delta)^{H-1-t}u_t.
\]
If $S^->0$, every compatible path has
\begin{align*}
 \log(1+\ell_H/C_H)-\gamma\log(S^+/S_{\rm ref})
 &\leq M_H\\
 &\leq\log(1+u_H/C_H)-\gamma\log(S^-/S_{\rm ref}).
\end{align*}
The bounds are sharp in the accounting-only class in which all interval combinations are feasible. With production, financing, or resource restrictions they remain valid outer bounds, but need not be sharp.
\end{proposition}
\begin{proof}
Iteration of the stock recursion gives the displayed affine stock formula. Its coefficients are nonnegative, so stock lies in $[S^-,S^+]$. The discrepancy increases with current investment $J_H$ and decreases with stock $S_H$ when $\gamma\geq0$. Substituting the appropriate bounds proves the inequalities. In the accounting-only interval class, set all earlier investments to their upper bounds and current investment to its lower bound for the lower endpoint; reverse the choices for the upper endpoint. These histories attain both endpoints. Additional feasibility restrictions may exclude them, which proves the stated limitation.
\end{proof}

\section{When instant adoption is physically unavailable}
Consider a required complementary stock $\kappa>0$ for full adoption. Full adoption at date $H$ is possible only if $S_H\geq\kappa$. Suppose investment capacity is bounded by an exogenous $\bar J\geq0$ each date, with no borrowing of future organizational capacity. The capacity model's complete feasible investment set is $0\leq J_t\leq\bar J$; a specified outside resource endowment finances these inputs. Stronger resource constraints are considered immediately below.

\begin{theorem}[A sharp capacity-only acceleration bound]\label{reach}
In the capacity-only model, full adoption at $H$ is feasible if and only if
\[
 \kappa\leq(1-\delta)^HS_0+
                  \bar J\sum_{t=0}^{H-1}(1-\delta)^t.
\]
In particular, date-zero full adoption is infeasible when $S_0<\kappa$. For $\delta=0$ and $\bar J>0$, the earliest feasible date is
$\max\{0,\lceil(\kappa-S_0)/\bar J\rceil\}$.
\end{theorem}
\begin{proof}
The stock recursion expresses $S_H$ as its inherited stock plus a nonnegative weighted sum of earlier investments. Replacing every investment by its capacity gives the largest attainable stock, which is the right side. This proves necessity; choosing every investment at capacity attains that stock and proves sufficiency. At $H=0$ no new investment has matured. If $\delta=0$, the maximum stock is $S_0+H\bar J$, so the smallest nonnegative integer satisfying its inequality is the stated ceiling.
\end{proof}

When investment must also satisfy $J_t+C_t\leq F_t$ for available real output $F_t$, the capacity bound is only necessary. It is not sufficient unless the path also leaves nonnegative consumption and obeys financing constraints. For example, take $S_0=1$, $\delta=0$, $\bar J=1$, $H=1$, $\kappa=2$. Capacity allows full adoption, but an endowment $F_0=1$ with a subsistence requirement $C_0\geq1$ forces $J_0=0$, so adoption remains infeasible. This example separates implementation capacity from actual resource availability.

\section{A feasible accelerated policy can lower the horizon residual}
We now specify both sectors and their resource use. A firm has one unit of labor, a stock $S=2$, and no physical capital. A policy chooses adoption $d\in\{0,1\}$ with the requirement $S\geq2d$. Adoption requires $d/2$ units of labor for own-account organizational investment. Thus
\[
 \ell=d/2,\quad J=\ell,\quad C=\eta(d)S(1-\ell),
 \quad Y=C+J,\quad S'=S+J,
\]
where $\eta(0)=1$ and $\eta(1)>1$. Labor is exactly split between final output and organizational investment; own-account output is valued at one per unit. Both policies have nonnegative consumption and are financed from their own production. Future capital exists and is not counted as current consumption. The adopted policy is a feasible, costly deployment comparison, not assumed to be the firm's unconstrained optimum.

Use the full-input residual $A=Y/(SL)$ with $L=1$ and the expensed measured residual $\widetilde A=C/L$. The actual benchmark is $d=0$; accelerated deployment is $d=1$.

\begin{proposition}[An explicit negative diffusion difference]\label{negative}
In this economy,
\[
 A^{\rm actual}=1,\qquad A^{\rm accelerated}=\tfrac12\eta(1)+\tfrac14.
\]
Consequently $D=\log A^{\rm accelerated}-\log A^{\rm actual}<0$ whenever $1<\eta(1)<3/2$, despite an improvement in final-sector efficiency. For $\eta(1)=11/10$, the full residual is $4/5$ under acceleration.
\end{proposition}
\begin{proof}
Without adoption, $\ell=J=0$, so $C=Y=2$ and $A=1$. Under adoption $\ell=J=1/2$, $C=\eta(1)$, and $Y=\eta(1)+1/2$. Dividing by $S=2$ gives the formula. It is below one precisely when $\eta(1)<3/2$. At $11/10$ it equals $11/20+1/4=4/5$. All labor, stock, and consumption constraints were explicitly verified by the displayed technology and choices.
\end{proof}

The negative residual is a sectoral-allocation effect in a declared aggregate input index. It does not say the underlying final-sector efficiency fell. If instead ``true productivity'' were defined to be the parameter $\eta(d)$ itself, the sign would be positive by definition; that would be a different target. The example shows why an accelerated policy, its resource costs, and the residual inventory must be fixed before interpreting $D_H$.

\begin{proposition}[Actual-path data do not identify acceleration]
Even complete observation of the actual path's $C,J,S,L,d$ does not identify $D$ in this model when $\eta(1)$ is unobserved. It can identify $M$ under the known accounting convention.
\end{proposition}
\begin{proof}
Compare two economies identical at $d=0$ but with $\eta(1)=6/5$ and $\eta(1)=2$. All actual quantities and residuals coincide. Both accelerated paths are feasible with the same labor cost and initial stock, but their accelerated residuals are $17/20$ and $5/4$, yielding different signs for $D$. Proposition~\ref{accounting} computes actual-path $M$ from its observed quantities in both models. Therefore knowledge of $M$ does not fill the missing adoption response needed for $D$.
\end{proof}

\section{Attribution and the remaining empirical gap}
Once both paths are defined, the overall residual gap satisfies the identity
\[
 \log A_H^{\rm accelerated}-\log\widetilde A_H^{\rm actual}=D_H+M_H.
\]
The same intangible investment can affect the feasible real adoption path and the accounting discrepancy without being counted twice in this identity: $D_H$ compares full residuals across two feasible paths, while $M_H$ changes the accounting convention on the actual path only. Calling both quantities separate dollar losses and adding them to welfare would have no such justification.

Identifying the full target still requires an investment measurement model, initial stock and depreciation, a complete input inventory, and credible variation in deployment costs or structural restrictions identifying the feasible accelerated path. Later measured realization is a useful test of those restrictions, not a logical consequence of current undermeasurement. Financing, firm optimization, prices, aggregation across firms, and policy equilibrium remain unsolved here; hence the original target remains unresolved.

\begin{thebibliography}{9}
\bibitem{ai} E. Brynjolfsson, A. Korinek, and A. Agrawal (2024), \emph{The Economics of Transformative AI: A Research Agenda}. Section 3.1, Economic Growth, printed p. 10, adoption/diffusion paragraph inspected 9 October 2026. PDF accessible; no claim that it supplies this intangible-accounting model. \url{https://digitaleconomy.stanford.edu/app/uploads/2024/11/ETAI-White-Paper.pdf}.
\end{thebibliography}

\end{document}
